\section{Engaging Citizens with Open Government Data: The Value of Dashboards Compared to Individual Visualizations}
Ich habe den Artikel \textit{„Engaging Citizens with Open Government Data: The Value of Dashboards Compared to Individual Visualizations"} \cite{simonofski2022} gewählt. Der wissenschaftliche Artikel ist öffentlich zugänglich und wurde von ACM Digital Government: Research and Practice (2022) veröffentlicht.

\smallskip
\noindent
Der Artikel untersucht, wie Bürger:innen besser mit Open Government Data (OGD) eingebunden werden können. Die Autoren identifizierten 16 Designprinzipien für OGD-Dashboards und evaluierten diese mit 108 Teilnehmenden. Im Folgenden fasse ich die wichtigsten Erkenntnisse zusammen und gehe darauf ein, wie sie sich auf unser Projekt beziehen.

\begin{paracol}{2}
%
\smallskip
\begin{itemize}
    \item \textbf{Dashboards erhöhen das Verständnis}

    Gut gestaltete Dashboards helfen Bürger:innen, OGD deutlich besser zu verstehen als einzelne, isolierte Visualisierungen.

    $\rightarrow$ Für unsere App bedeutet das: Wir sollten Daten nicht isoliert anzeigen, sondern in einem strukturierten, zusammenhängenden Kontext darstellen.
    %
    \item \textbf{Heterogene Nutzergruppen berücksichtigen}

    Die Studie zeigt, dass OGD-Nutzer:innen sehr unterschiedliche technische Hintergründe haben -- von erfahrenen Expert:innen bis hin zu Gelegenheitsnutzer:innen ohne Vorkenntnisse.

    $\rightarrow$ Das spiegelt sich direkt in unseren Personas wider: Von Benjamin als technikaffinern Studierenden bis zu Luca als Gelegenheitsnutzer muss unsere App für alle verständlich sein.
    %
    \item \textbf{Klare Beschriftungen und Erklärungen}

    Ein häufiges Problem bei OGD-Plattformen ist fehlendes Kontextwissen: Nutzer:innen verstehen nicht, was die Daten bedeuten oder woher sie stammen.

    $\rightarrow$ Unsere App muss jeden Datensatz mit verständlichen Beschriftungen und Quellenangaben versehen, damit auch nicht-technische Nutzer:innen die Daten einordnen können.
    %
\switchcolumn
    %
    \item \textbf{Geeignete Visualisierungsformen wählen}

    Nicht jede Visualisierungsform eignet sich für jeden Datensatz. Die Studie betont, dass die Wahl der richtigen Darstellungsform entscheidend für das Verständnis ist.

    $\rightarrow$ Da wir eine breite Bandbreite an unterschiedlichen Wiener Open Datensätzen visualisieren, müssen wir für jeden Datensatz die passende Visualisierungsform wählen.
    %
    \item \textbf{Datenqualität transparent machen}

    Bürger:innen müssen darauf vertrauen können, dass die angezeigten Daten korrekt und aktuell sind. Fehlende Transparenz über die Datenqualität ist ein zentrales Hindernis für die Nutzung von OGD.

    $\rightarrow$ Besonders relevant im Hinblick auf unsere negative Persona: Unsere App muss sicherstellen, dass Daten aus vertrauenswürdigen Quellen stammen und nicht für manipulative Zwecke missbraucht werden können.
\end{itemize}
%
\end{paracol}
